\section{Gestión de datos}
\label{sec:sd5-5-2}

Esta sección desarrolla RT-05.02–10: motor y paradigma, consistencia y disponibilidad, garantías de escritura, autoridad offline, gobierno de maestros, calidad, separación analítica, retención y exportación. Cada decisión se aplica a los dominios de 5.1 y define su comportamiento ante fallos.

\subsection{Paradigma y motor de persistencia por dominio}
\label{sec:sd5-5-2-1}

CAP designa consistencia, disponibilidad y tolerancia a particiones. Durante una partición no se promete consistencia global inmediata junto con disponibilidad de todas las escrituras: la autoridad de cada agregado mantiene sus invariantes y las proyecciones remotas pueden quedar desactualizadas. Los compromisos por dominio se desarrollan a continuación y en 5.2.3.

Los quince nombres lógicos de SD4 separan responsabilidades de escritura, sin exigir quince motores. El Anexo 5-B integra dominio, ubicación, motor y autoridad. PostgreSQL 16/PostGIS sostiene relaciones, restricciones e historia; Aurora concentra los dominios centrales y los PostgreSQL de sitio sostienen inventario y preparación bajo una sola autoridad por agregado. BD\_RUTAS tiene autoridad central en Aurora y copias de consulta autorizadas en sitio o móvil.

Frente a una partición de red, se conserva consistencia dentro del perímetro de la autoridad: bodega puede confirmar contra su PostgreSQL local y terreno puede registrar hechos autorizados durablemente. La consolidación central queda pendiente. Una copia indicativa de stock no permite crear reservas. DynamoDB acepta ingesta térmica eventual, Amazon S3 en formato Parquet conserva detalle y Redshift explota series y agregados. Amazon ElastiCache (Valkey 8.1) contiene proyecciones reconstruibles y S3\_DOCS objetos versionados. Las credenciales de AWS son administración, no un motor de persistencia adicional.

La elección relacional evita saldos y cobros divergentes; el modelo columnar reduce barridos analíticos sin competir con OLTP. La disponibilidad local no concede autoridad para modificar compromisos centrales ni para crear permisos offline.

\subsection{Transaccionalidad, aislamiento e idempotencia}
\label{sec:sd5-5-2-2}

La telemetría usa productor, arranque y secuencia para procesar una sola vez cada mensaje de red móvil o satélite; las lecturas del mismo mensaje conservan su sensor. M9 detecta huecos y solicita el tramo faltante, ubica por secuencia las llegadas desordenadas y siempre procesa una crítica tardía: retiene los SSCC afectados o abre la evaluación de retiro si ya fueron entregados. Una recuperación tardía nunca cierra una crítica ni una retención. La serie fina reemplaza el resumen satelital para formar la evidencia.

Negocio, auditoría y outbox se escriben en la misma transacción del módulo dueño. Bajo lectura confirmada, las operaciones de saldo/reserva/cobro bloquean su agregado y validan invariantes; una clave única protege idempotencia, sin depender de consultar antes de insertar. El shipper/trabajador de SD4 publica la outbox, el motor no envía directamente a la cola.

Inbox y efecto se confirman en una transacción: caída previa revierte ambos; posterior al commit devuelve el resultado, sin repetir efectos. Las proyecciones no retroceden versión, pero hechos aditivos/sanitarios/financieros tardíos se conservan. Un gap reordena y solicita faltantes; tras agotar reintentos escala y concilia. Actor, dispositivo, correlación, esquema/versión e instantes de ocurrencia/registro permiten reconstrucción RT-05.03/.19.

Offline, Room/SQLite persiste hecho e identidad/evidencia antes del acuse \ensuremath{\leq}2 s. El móvil verifica autorización preemitida por Keycloak, puede cerrar localmente la entrega autorizada con consolidación pendiente. La carga/hash central confirma después el objeto. La aceptación local no afirma que la nube ya conoce el hecho ni concede autoridad para reservar contra una copia de consulta.

\subsection{Autoridad de escritura y comportamiento en partición}
\label{sec:sd5-5-2-3}

Cada agregado conserva un escritor autorizado. Sin WAN, bodega consulta y escribe contra PostgreSQL local, una reserva allí validada está confirmada localmente, aunque su consolidación central siga pendiente. DMS replica exclusivamente Talca hacia un esquema central lector separado. Las demás sedes publican eventos; los consumidores centrales deduplican efectos de negocio. CDC no vuelve a aplicar esos efectos ni sustituye la outbox. BD\_RUTAS se administra centralmente y los itinerarios/manifiestos descargados permiten ejecutar la ruta autorizada offline.

Sin WAN, el móvil registra intentos, entregas, cobros permitidos y evidencias bajo permisos preemitidos; stock, precio y crédito descargados son indicativos para decisiones centrales. La incertidumbre POS, autorización segregada de pago alternativo y conciliación siguen 5.1.8: un timeout no prueba rechazo ni garantiza ausencia de doble cargo.

La Tabla 5.1 compara los límites declarados en SD4 y su consecuencia observable.

\parte{tablas/limites_offline}

Los plazos de la Tabla 5.1 corresponden a objetos con funciones distintas. La copia de identidad permite consultar datos durante 24 horas y la credencial autoriza las operaciones del turno durante su propia vigencia. El manifiesto de 26 horas permite verificar los relevos inscritos previamente, pero no prolonga una credencial vencida. En bodega, los relevos previstos a las 8 y 16 horas permiten cubrir la autonomía de 24 horas con credenciales de turno válidas. Keycloak firma y renueva el manifiesto mientras existe conexión. Si el enlace se pierde, el dispositivo comprueba localmente la firma y la vigencia. La evidencia de entrega queda guardada de forma durable antes de confirmar el cierre local autorizado, y se carga al centro cuando vuelve la conexión. Para una ruta de hasta 14,5 horas, la planificación debe prever un punto de enlace para renovar la credencial antes de cumplir las 14 horas sin cobertura. Si el enlace falla y la credencial vence, el terminal conserva los registros ya capturados, bloquea nuevas operaciones y activa la contingencia de despacho hasta recuperar autorización. El Anexo 5-J comprueba ambos casos.

\subsection{Gestión de datos maestros}
\label{sec:sd5-5-2-4}

El dueño funcional propone cambios, un rol habilitado aprueba y el módulo de maestros publica versión y vigencia sin solapamiento. Leandro Chamorro coordina datos según SD1; responsables del CLIENTE validan Comercial, Finanzas, bodega y Calidad. Un conflicto de RUT/GTIN/código externo se aísla con original y decisión; no se fusionan lotes por semejanza ni se activan maestros incompletos por defecto.

La capa anticorrupción de SD4 traduce ERP/WMS/cadenas a las claves y unidades canónicas. Contratos OpenAPI/AsyncAPI, autenticación entre sistemas, modo, volumen, ventana y fallo se remiten a los Anexos 4-G, 4-H y 4-I de SD4; las versiones y el preaviso mínimo de seis meses de RT-05.17 conservan lectores de la ola previa. La publicación por eventos invalida copias y deja versión/instante visible; un sitio sin enlace no crea un maestro central alternativo.

\subsection{Calidad de datos}
\label{sec:sd5-5-2-5}

La solución valida los datos al capturarlos, integrarlos y utilizarlos. Cada validación conserva el dato original y registra la regla aplicada, su versión, la causa, el responsable y la resolución. En la línea base del caso, el 41 \% de las recepciones del producto retirado no registraba el lote. Para cerrar una recepción, deben estar completas el 100 \% de las líneas de productos que exigen lote. El identificador SSCC se exige cuando la unidad logística alcanza el hito que lo requiere. Su ausencia es válida en los estados \texttt{VACIA} o \texttt{ABIERTA}, y para los productos exentos.

ISO/IEC 25012 (International Organization for Standardization, 2008) orienta la evaluación de la completitud, exactitud, consistencia, unicidad, validez, oportunidad y credibilidad de los datos. El Anexo 5-C define los valores admitidos y las causas de rechazo. El Anexo 5-H documenta las fórmulas y el tablero de calidad, desde el que se puede consultar el registro y la decisión asociados. El linaje de RT-05.10 se documenta conforme al Formulario T-12.

La desviación monetaria se calcula respecto del valor contado en el sitio y se evalúa por separado de la exactitud por cantidades. LafroX acepta una desviación de hasta el 0,3 \% y exige explicar cada diferencia. El 2,3 \% corresponde a la línea base del caso. Las lecturas térmicas inválidas se conservan marcadas y la pérdida de cobertura se mide sin eliminar la evidencia del indicador.

\subsection{Separación transaccional-analítica y modelo de explotación}
\label{sec:sd5-5-2-6}

El indicador de entrega completa y a tiempo (OTIF) se calcula como 100 \ensuremath{\times} el cociente entre los pedidos comprometidos que se entregaron completos dentro de la ventana \texttt{ORIGINAL} y el total de pedidos comprometidos en el período. Cada pedido se cuenta una sola vez y el denominador incluye los pedidos no entregados. La promesa se identifica por el tipo \texttt{ORIGINAL}, sin filtrar por vigencia, de modo que las revisiones posteriores no cambian la base de comparación. El SD3 fija metas de 90 \% en el mes 15, 93 \% en el mes 19 y 95 \% en el mes 32. Si el denominador es cero, el indicador se informa como no aplicable.

La exactitud por cantidad compara el registro del sistema con el conteo físico en unidades homologadas, para cada combinación de sitio, ubicación, producto, lote y unidad de medida. Su fórmula es 100 \ensuremath{\times} max(0, 1 \ensuremath{-} \ensuremath{\Sigma}|cantidad sistema \ensuremath{-} cantidad conteo| / \ensuremath{\Sigma}|cantidad conteo|). Si tanto la base física como la del sistema son cero, se informa que no hay base de cálculo. Si el conteo es cero y el sistema registra una cantidad distinta de cero, se muestra la diferencia y el porcentaje se informa como no aplicable.

La desviación monetaria se calcula como 100 \ensuremath{\times} \ensuremath{\Sigma}|valor sistema \ensuremath{-} valor físico| / valor contado del sitio. Los valores se expresan en CLP y se obtienen con el mismo criterio de valorización. Si la base contada es cero, se registra la decisión adoptada. El criterio de aceptación de LafroX es \ensuremath{\leq}0,3 \%, con independencia de la exactitud por cantidades.

El indicador de integridad de recepción solo considera las líneas de productos que requieren lote al cierre. La comprobación de SSCC se hace en el hito correspondiente de la unidad logística, sin penalizar los estados \texttt{VACIA} y \texttt{ABIERTA} ni los productos exentos.

Los tableros operativos consultan datos por API y proyecciones separadas, con un desfase máximo de cinco minutos. El cierre se publica como máximo dos horas después del retorno del último camión. Los tableros gerenciales, en Redshift y QuickSight, tienen una latencia máxima de cuatro horas. El cálculo del costo de servicio comienza en la Etapa 2.

El procesamiento analítico de mayor carga se ejecuta en Redshift y Amazon S3 mediante Glue, separado del procesamiento transaccional en línea (OLTP). Las consultas operativas puntuales muestran la marca de actualización de sus datos. En Talca, la captura de cambios (CDC) permanece en un esquema de lectura separado para evitar duplicar el efecto de los eventos de negocio.

La ampliación mantiene las exclusiones del SD3 y SD4: linaje automatizado, portal de desarrolladores y las capacidades predictivas deseables. La estimación determinista de la vida útil remanente sí se incluye, según el SD13, sección 13.3.2. Utiliza parámetros cinéticos documentados, las series térmicas y la asociación entre sensor y lote descrita en sección 5.1.7, sin usar modelos de aprendizaje. Se registra una sola vez como capacidad analítica de RT-05.30. INN-03 desarrolla su ejecución distribuida y su uso en las decisiones operativas (Subdocumento 13, sección 13.3.2).

El CLIENTE puede crear informes desde las consolas mediante la autoría integrada de QuickSight, con permisos de lectura y creación diferenciados. El modelo semántico de sección 5.1.11 y el Anexo 5-H definen el nivel de detalle, las relaciones, las fórmulas y los casos sin base. El catálogo indica audiencia, etapa y latencia, y permite filtrar por período, sitio, producto y canal, consultar el hecho y el documento de origen, exportar los datos y consultar su linaje documental. Los informes admiten filtros guardados y envíos programados bajo los mismos permisos, sin intervención del ADJUDICATARIO.

Cada porcentaje muestra su base e indica cómo se trata un denominador cero. Las cantidades solo se comparan tras homologar sus unidades, y los pedidos no entregados se mantienen en la base del OTIF. La marca de consolidación y los registros pendientes permiten distinguir los datos actualizados de los que aún están por consolidar.

\subsection{Retención, archivado, eliminación y protección}
\label{sec:sd5-5-2-7}

La temperatura y los archivos originales del termógrafo y del sensor independiente se conservan cinco años en Amazon S3 con bloqueo de objetos en modo cumplimiento. \texttt{tel\_\allowbreak{}archivo\_\allowbreak{}termino} incluye ambos registradores; se verifica integridad, el TC58e firma el hash y la memoria solo se libera tras confirmación de nube. El registro térmico no almacena GPS; la historia del SSCC se determina por compartimento, intervalos y escaneos.

El capítulo 15 del caso identifica como RT-05.10 la retención histórica y de auditoría, en las Bases Técnicas Transversales, RT-05.10 corresponde al catálogo con linaje automatizado y RT-16.10 a retención de auditoría. Se solicita aclarar al CLIENTE la correspondencia de códigos, sin presumir una respuesta oficial. Mientras se aclara, se aplican íntegramente los plazos del caso por clase de dato y las obligaciones transversales de retención; el catálogo automatizado deseable mantiene su alcance no ofertado. La discrepancia de identificador no suspende ni reduce la conservación (Caso, cap. 15, p. 26,Bases Técnicas Transversales, caps. 5 y 16, pp. 12 y 29).

La conservación depende de la clase de información y del expediente que deba interpretarse. Los documentos y las evidencias de entrega se conservan seis años. El expediente sanitario permanece durante el mayor plazo entre la vida útil del producto más seis meses y cinco años. El detalle térmico se conserva cinco años, la geolocalización personal doce meses en total y la auditoría siete años. Los registros técnicos permanecen doce meses disponibles en línea y veinticuatro meses adicionales en archivo. Los respaldos se retienen treinta y cinco días, según el SD4. El Anexo 5-D detalla los plazos por dominio y atributo, y el Anexo 5-E define su clasificación y cifrado. Guardar un objeto en Amazon S3 no determina por sí solo su plazo de conservación.

Maestros, reglas y calibraciones se conservan mientras interpreten un expediente vigente. Un rango temporal no autoriza eliminar toda una partición si incluye expedientes abiertos. La purga verifica filas, versiones, derivados, cachés, réplicas y copias legibles, una retención legal/Object Lock activa aplaza la eliminación elegible y deja evidencia. Eliminar coordenada requiere impedir reconstrucción desde auditoría; geocerca de configuración no es posición personal. No se destruye clave compartida de evidencia vigente. El archivo exige acceso auditado y restauración probada; la ventana histórica de migración no es plazo de borrado.

\subsection{Exportación e intercambio masivo}
\label{sec:sd5-5-2-8}

Se propone una exportación íntegra, autónoma y sin cobro adicional, operada por un rol autorizado del CLIENTE. Incluye datos aceptados e historia, maestros y equivalencias, reglas y calibraciones, auditoría, esquemas/diccionario, manifiestos y objetos con sus referencias; respeta los permisos y los plazos de retención. Se entrega CSV UTF-8 o JSON para relaciones, Parquet para series y originales de objetos, con esquemas abiertos y sin exigir software propietario para recuperar su contenido.

El trabajo registra snapshot por autoridad, instante de corte y marcas de agua; captura los deltas hasta una frontera común conciliada y publica la relación de autoridades rezagadas. Una exportación incompleta se declara parcial y no se ofrece como copia íntegra. El manifiesto contiene versión de esquema, alcance, conteos, sumas de control, hashes SHA-256 y referencias de objetos; la validación compara la descarga con ese manifiesto.

El CLIENTE inicia, consulta avance y descarga por enlace autenticado/temporal. Se auditan solicitante, alcance y accesos. Como parámetros iniciales de diseño, el enlace vence en 24 h y puede renovarse sin repetir paquetes; estos se purgan a siete días conservando auditoría. Identidad de trabajo y bloques verificados permiten reanudación sin duplicar/omitir datos. La aceptación ensayará reconstrucción de datos y apertura de objetos fuera de la plataforma.

La exportación de gran volumen se ejecutará como trabajo asíncrono y mantendrá disponible la sesión para las demás operaciones autorizadas. Al completarse, el servicio de notificaciones existente avisará a la persona solicitante y ofrecerá acceso a la descarga bajo sus permisos (RT-16.29).

La carga masiva operativa de RT-05.22 usa consolas y trabajos por lotes de SD4, Anexo 4-C. Un rol CLIENTE autorizado carga CSV UTF-8/JSON según contrato versionado del módulo. Se validan formato, permisos, claves, dominios y referencias antes de aplicar, conservando identidad de origen/operación. Se aceptan registros válidos independientes; dependencias/agregados inseparables se rechazan completos. El informe consultable/descargable contiene estado, conteos y fila/clave, causa y regla de cada rechazo. Corregir y reenviar rechazados conserva identidad sin repetir aceptados, negocio, auditoría y acuse siguen 5.2.2. Descarga y exportación programada mantienen permisos, formatos y manifiesto.
